\chapter{الإجراءات الداخلية}
\label{ch:builder-internal-sops}

\section{ماذا ستبني في هذا الفصل؟}
\label{sec:builder-internal-sops-intro}

الإجراء الداخلي مسار عمل من خطوات متتابعة يبدأ عند حدث تحدده، مثل إنشاء سجل أو إرسال نموذج. وقد تطلب خطواته موافقة، أو تُنشئ مهمة، أو ترسل إشعاراً، فلا يعتمد إنجاز العمل على أن يتذكر أحد الخطوة التالية. ويشرح هذا الفصل الأداة كاملة: الصفحة الرئيسية، والمحرر، والمشغّلات وأنواع الخطوات، وإعداد كل خطوة ومعالجة أخطائها، ثم الاختبار والنشر ومتابعة عمليات التشغيل والرد على الموافقات.

وتُحفظ تعديلاتك في مسودة لا تعمل حتى تنشرها، لذلك يمر كل إجراء بالدورة نفسها: تبنيه، ثم تختبره، ثم تنشره، ثم تتابع تشغيله.

\begin{rcnote}[مثال: شركة الأفق]
عند إنشاء طلب زيارة جديد يطلب إجراء داخلي موافقة سارة الخطيب، وبعد موافقتها يُنشئ مهمة تحضير الزيارة لعمر حداد.
\end{rcnote}

\begin{figure}[H]
\centering
\begin{tikzpicture}[node distance=0.5cm]
\node[stepBox] (a) {\RL{المشغّل: ما الذي يبدأه؟}};
\node[stepBox,below=of a] (b) {\RL{الخطوات وإعدادها}};
\node[stepBox,below=of b] (c) {\RL{اختبار: محاكاة آمنة}};
\node[stepBox,below=of c] (d) {\RL{مراجعة ونشر}};
\node[stepBox,below=of d] (e) {\RL{متابعة عمليات التشغيل}};
\draw[flowsTo] (a) -- (b);
\draw[flowsTo] (b) -- (c);
\draw[flowsTo] (c) -- (d);
\draw[flowsTo] (d) -- (e);
\end{tikzpicture}
\rclegend
\caption{دورة حياة الإجراء الداخلي}
\label{fig:builder-internal-sops-overview}
\shotnote{يرتب المخطط صفحات هذا الفصل بحسب دورة الإجراء. فالاختبار يسبق النشر لأن المحاكاة لا تُحدث أي أثر، والمتابعة تأتي بعد النشر لأن الإجراء لا يبدأ أي تشغيل قبله.}
\end{figure}

\section{صفحات هذا الفصل}
\label{sec:builder-internal-sops-pages}

تبدأ من صفحة «الإجراءات الداخلية» في قسم «الأساسية» من أدوات البناء، أما بقية الصفحات فتشرح أجزاء المحرر وتبويبات الصفحة نفسها. واتبع ترتيبها في المرة الأولى، ثم عُد إلى أي صفحة عند الحاجة.

\begin{table}[H]
\centering
\caption{تشغيل العمل بالإجراءات الداخلية}
\label{tab:builder-pages-internal-sops}
\rcZebra
\begin{tabularx}{\linewidth}{R{4.2cm} Y}
\rcHeaderRow
\thh{الصفحة وقسمها} & \thh{ماذا تقدم هذه الصفحة؟}\\
الإجراءات الداخلية — الأساسية & تعرض الإجراءات في مجلدات مع حالتها وصحة تشغيلها، ومنها تنشئ إجراءً جديداً.\\
محرر الإجراء وأوضاعه — المحرر & رأس المحرر وأوضاع «بناء» و«اختبار» و«عمليات التشغيل»، وقائمة الخطوات واللوحة وأدواتها.\\
المشغّلات: ما الذي يبدأ الإجراء؟ — المحرر & مجموعات المشغّلات، ومشغّلات السجلات والنماذج والمهام والموافقات، وتغيير المشغّل.\\
خطوات التحكم في المسار — المحرر & التفرّع الشرطي والتكرار والتأخير والانتظار والإنهاء واستدعاء إجراء آخر.\\
خطوات السجلات والنماذج — المحرر & البحث عن السجلات وإنشاؤها وتحديثها، وطلب تعبئة النماذج.\\
خطوات المهام والموافقات والإشعارات والمستندات — المحرر & إنشاء المهام وطلب الموافقة والتوقيع وإرسال الإشعارات وإنشاء المستندات.\\
إعداد الخطوة: المدخلات والبيانات والتعبيرات — المحرر & لوحة إعداد الخطوة، وإدراج البيانات من الخطوات السابقة، والتعبيرات والاسم المرجعي.\\
معالجة الأخطاء وتخطي الخطوات — المحرر & الإيقاف أو مسار الفشل عند الخطأ، وإعادة المحاولة، وتخطي الخطوة.\\
الاتصالات والمتغيرات — الإجراءات الداخلية & حسابات الخدمات الخارجية، والقيم المشفرة التي تستخدمها التعبيرات.\\
اختبار الإجراء — المحرر & المحاكاة الآمنة والتشغيل الفعلي واختبار خطوة واحدة.\\
عمليات التشغيل وسجلها — المحرر والإجراءات الداخلية & حالات التشغيل، وقراءة النتائج على اللوحة، وإعادة التشغيل أو إلغاؤه.\\
الرد على طلبات الموافقة — طلبات بانتظارك & كيف يرد المعتمد، وما يحدث بعد كل رد.\\
المراجعة والنشر وسجل الإصدارات — المحرر & ما يجب إصلاحه والتغييرات قبل النشر، والإصدارات المقفلة والإيقاف المؤقت.\\
\end{tabularx}
\end{table}
